% (c) Danny Lee Eldridge · AEGIS Community License · CC BY-NC-SA 4.0
\documentclass[aps,prd,twocolumn,superscriptaddress,nofootinbib,showkeys]{revtex4-2}
\usepackage{amsmath,amssymb,hyperref,booktabs}
\begin{document}
\title{The Quark Thread: A Twin Slip-Ladder of Supergravity Minima for the Quark Sector of a Modular Weave Model, Tested Step by Step}
\author{Danny Eldridge}
\email{dannyeldridge54@gmail.com}
\affiliation{Independent researcher, Oklahoma City, Oklahoma, USA}
\date{20 November 2026}
\begin{abstract}
In the ``Weave'' model the inflaton is a modular field that ends inflation at the fixed point $\rho=e^{i\pi/3}$, and the lepton moduli sit at the fixed points $\rho$ and $i$~\cite{P1,P2}. We ask where the quarks fit. Thirteen pre-registered tests (rules committed with a checksum before computing) give a mixed but sharp picture.
(i) Neither the inflaton's endpoint nor the late-time dark-energy field can be the quark modulus: the first is $\sim\!1600\times$ too close to $\rho$, the second would make quark masses drift today.
(ii) The lepton Koide relation does not extend to quarks ($Q_u=0.887$, $Q_d=0.742$).
(iii) A third, heavy modulus $\tau_q$ ($m\gtrsim2\times10^{10}$\,GeV) is compatible with the early universe, and anomaly-free modular weights give $\bar\theta=0$ without an axion~\cite{FST}.
(iv) The supergravity potentials $W\propto(j-1728)^{m/2}/\eta^6$ have CP-violating minima just off $\rho$; we compute them directly up to $m=32$ and find $|\epsilon_m|=|\tau-\rho|\simeq0.0251/(m+\tfrac1{16})$ at a fixed phase $\simeq50^\circ$ (agreeing with the analytic form of Ref.~\cite{NPP} to $\le1.3\%$, independently confirmed to $<10^{-5}$). Mixed classes $(m,n)$ sit on the CP-conserving arc $|\tau|=1$ at a position set by $m{:}n$ alone.
(v) Splitting the quark thread, down quarks on rung $m=1$ and up quarks on a higher rung, removes the large coupling the single-modulus fits need. In the weight-8/6/4 model of Ref.~\cite{PT} it lowers $\chi^2$ from 20.4 to 5.8 with all couplings $\mathcal O(1)$, but still misses $m_s/m_b$ ($0.54\times$), $|V_{cb}|$ ($1.27\times$) and $|V_{ub}|$ ($0.73\times$): a FAIL by the locked tolerances.
(vi) With $\mathcal O(1)$ coefficient spread, the up-quark rung is only bounded to $5$--$14$ (68\%).
The quark sector is therefore hosted and consistent, not predicted. We freeze four predictions for future data.
\end{abstract}
\keywords{modular flavour symmetry; quark masses; CKM; strong CP; supergravity; modulus stabilisation; pre-registration}
\maketitle

\section{Setting}
The Weave model~\cite{P1} drives inflation by the string coupling $\tau_1$ on a species-scale plateau and ends it at $\rho$, where a residual $Z_3$ survives. The charged-lepton matrix is then a $Z_3$ circulant, $\sqrt{M_\ell}=a+bC+\bar bC^2$, which gives the Koide form; the neutrino modulus sits at $i$~\cite{P2}. Fixed points conserve CP, and quark moduli placed there give a Jarlskog invariant $J\le5\times10^{-15}$, far below $3\times10^{-5}$~\cite{P2}. The quark sector therefore needs its own modulus $\tau_q$ away from the CP-symmetric lines. We test, step by step, where $\tau_q$ can sit and what it predicts. All steps were locked in advance (\texttt{field\_tests/rules/PREREG\_Q*}); deviations were written down before any rerun; results are in Table~\ref{tab:steps}.

\section{Which field carries the quarks}
\textit{Q1, inflaton endpoint.} Near $\rho$, level-3 $A_4$ quark models need $\tau=\rho+\epsilon$ with $|\epsilon|\approx0.03$~\cite{PT}. The inflaton's leftover canonical offset after reheating is $\le10^{-5}M_P$, i.e.\ $|\epsilon|\le1.2\times10^{-5}$ (Kähler weight $h=1$): $\sim1600\times$ too small. \textit{The dark thread} (the late dark-energy field) rests at $i$, a CP-conserving point, and is still rolling; tying quark masses to it would make them drift by order one over gigayears, excluded by atomic clocks and equivalence-principle tests (our bound $\zeta<3.6\times10^{-6}$). Both are rejected.
\textit{Q2, Koide.} With masses run to $M_Z$~\cite{MZ}, $Q=\sum m/(\sum\sqrt m)^2$ is $0.8870\pm0.0010$ (up) and $0.7422\pm0.0025$ (down), far from $2/3$; the Brannen phases ($0.052$, $0.104$) are far from the lepton value $2/9$. The same code reproduces the leptons ($Q=0.6666645$, $\delta=0.222225$).
\textit{Q3, predictivity.} The anomaly-free modular quark models of Ref.~\cite{FST} fit the data with 15--22 real parameters for 8--10 observables; one variant cannot fit at all. What survives is structural: $\bar\theta=0$ for any $\tau_q$, with CP broken only by $\langle\tau_q\rangle$.
\textit{Q4, cosmology.} A displaced modulus decaying gravitationally dilutes the universe by $D=\tfrac16T_R(\phi_0/M_P)^2/T_{\rm dec}$ and shifts $N_*$ by $-\tfrac13\ln D$. With $T_R=1.34\times10^9$\,GeV and $H_{\rm inf}=5.2\times10^{12}$\,GeV, requiring $D\le10^3$, $T_{\rm dec}\ge5$\,MeV and $|\Delta n_s|\le0.003$ gives $m_{\tau_q}=1.7\times10^{10}$--$5\times10^{12}$\,GeV for $\phi_0=M_P$ ($8\times10^8$\,GeV upward for $\phi_0=0.1M_P$), uncertain by a factor $\sim2$.

\section{The supergravity minima}
With $K=-3\ln(2\,{\rm Im}\,\tau)$ and $W=H(\tau)/\eta^6$, $H=(j-1728)^{m/2}j^{n/3}$~\cite{NPP},
\begin{equation}
V=\frac{|H|^2}{(2y)^3|\eta|^{12}}\Big[\tfrac{4y^2}{3}\Big|\tfrac{H'}{H}-\tfrac{i\pi}{2}E_2+\tfrac{3i}{2y}\Big|^2-3\Big].
\end{equation}
We minimise $V$ in the fundamental domain (q-series to $q^{60}$, multi-start). The code reproduces the three minima of Ref.~\cite{NPP} ($m=1,2,3$) to $|\Delta\tau|\le5\times10^{-4}$. For $n=0$ and $m=1\ldots32$ the minima obey
\begin{equation}
\epsilon_m\equiv\tau_m-\rho\simeq\frac{0.0251\,e^{i\,50^\circ}}{m+1/16},
\label{eq:ladder}
\end{equation}
with the computed $|\epsilon_m|$ within $0.3$--$1.3\%$ of Eq.~(\ref{eq:ladder}) ($0.02395$ at $m=1$, $0.00313$ at $8$, $0.00079$ at $32$). The phase is stable at $49.5$--$50.9^\circ$ up to $m=14$; beyond, $|\epsilon|\sim10^{-3}$ and the optimiser cannot resolve the angle. An independent stationary-point solver (mpmath, root finding on $\nabla V$) agrees to $4\times10^{-7}$ ($m=1$), $1.6\times10^{-6}$ ($m=8$). Mixed classes land exactly on the arc $|\tau|=1$ (CP-conserving): $(1,1)$, $(2,2)$, $(3,3)$ all at $\tau\simeq-0.2378+0.9713i$, moving toward $\rho$ as $m$ grows and toward $i$ as $n$ grows. The minima thus form a ``twin ladder'': CP-violating rungs climbing to $\rho$, minima at $i$, joined by a CP-conserving arc. Eq.~(\ref{eq:ladder}) with no free parameter ($m=1$: $|\epsilon|=0.024$) lands below the band $0.028$--$0.038$ that single-modulus fits need~\cite{PT} (Q6, FAIL).

\section{A split quark thread}
Near $\rho$, with $\mathcal O(1)$ couplings, a family on rung $m$ has $m_2/m_3\simeq\kappa_2\epsilon_m$, $m_1/m_3\simeq\kappa_1\epsilon_m^2$~\cite{PT}. The down sector ($m_s/m_b=0.0199$, $\sqrt{m_d/m_b}=0.0317$) matches rung 1; the up sector ($0.0037$, $0.0028$) matches rungs 3--16 (Q7). Charged leptons would also match rung 1 by masses alone, but a rung-1 slip breaks the Koide relation: with random $\mathcal O(1)$ non-circulant terms of size $|\epsilon|$, $|\Delta Q|\approx0.37|\epsilon|$, i.e.\ $8.5\times10^{-3}$, versus a 3$\sigma$ tolerance $2.3\times10^{-5}$ (only $0.15\%$ of draws survive). Leptons therefore sit at the knot ($|\epsilon|\lesssim6\times10^{-5}$), and get their hierarchy from the circulant numbers, not from a slip (Q10).

\textit{CKM.} We built the mass matrices of Ref.~\cite{PT} (validated: rank one at $\tau=\rho$ with the paper's rows to $10^{-16}$), put $\tau_d$ on rung 1 and $\tau_u$ on rungs 8--10 (and 13, added after the fact because the two-moduli fit of Ref.~\cite{PT} has $|\epsilon_u|=0.0019$), and searched $\mathcal O(1)$ couplings ($[1/3,3]$, any phase). Tolerances were set in advance: mass ratios $\times/\div1.5$, $|V_{us}|$ $\pm10\%$, $|V_{cb}|$ $\pm20\%$, $|V_{ub}|$ $\pm30\%$, $J$ $\pm50\%$.
\begin{itemize}
\item Original model (Q8): one modulus $\chi^2=54$; split $\chi^2=39$ (rung 8). Up masses and $|V_{cb}|$ become right; $m_s/m_b$, $|V_{ub}|$ and $J$ are $\sim3\times$ low. FAIL.
\item Weight-8/6/4 model (Q12; chosen after Q8, as suggested in Ref.~\cite{PT}): one modulus $\chi^2=20.4$; split $\chi^2=5.8$--$6.0$ (rungs 8--10), $7.1$ (13). Passing: $m_d/m_b$, $m_c/m_t$, $m_u/m_t$, $|V_{us}|$ ($0.97\times$), $J$ ($0.91\times$). Failing: $m_s/m_b$ ($0.54\times$), $|V_{cb}|$ ($1.27\times$), $|V_{ub}|$ ($0.73\times$). FAIL, but the closest result, and without the coupling $f_u\approx30$ the single-modulus fit requires.
\end{itemize}
In both models the down sector wants a larger slip than rung 1, the outermost rung of the $(m,0)$ family. We did not widen the coupling box or allow non-integer $m$.

\textit{How sharp is the up rung?} In the equal-coefficient limit the measured ratio $g_u/g_d$ ($g=\sqrt{(m_2/m_3)\sqrt{m_1/m_3}}$) gives $m_u^*=8.04\pm0.11$ from precise $M_Z$ Yukawas~\cite{RUN}, rising to $9.2$, $9.5$, $9.8$ at $10^{9}$, $10^{12}$, $10^{16}$\,GeV; the down quarks stay on rung 1 at every scale. Including the $\mathcal O(1)$ spread ($\kappa$ log-uniform in $[1/3,3]$), the posterior is broad: rung 8 most probable ($8\%$), 68\% set 5--14, 95\% set 3--22. The robust statement is that up and down quarks sit on different rungs, the up rung roughly ten steps in.

\begin{table}[t]
\caption{\label{tab:steps}Locked steps (D1727--D1744 in the project log).}
\begin{ruledtabular}\begin{tabular}{lll}
step & question & result\\ \hline
Q1 & inflaton endpoint as $\tau_q$ & FAIL ($\times1600$)\\
Q2 & Koide for quarks & FAIL\\
Q3 & modular quark models predictive & fit only\\
Q4 & heavy $\tau_q$ vs cosmology & PASS ($\gtrsim2\times10^{10}$ GeV)\\
Q6 & zero-parameter $\epsilon$ vs fits & FAIL ($0.024$)\\
Q7 & split thread hosts both sectors & PASS (structure)\\
Q9 & two-anchor ladder map & map\\
Q10 & leptons on rung 1 & NO (Koide)\\
Q11 & computed minima, $m\le32$ & formula holds ($\le1.3\%$)\\
Q8 & CKM, original model, split & FAIL ($\chi^2$ 54$\to$39)\\
Q12 & CKM, weight-8/6/4, split & FAIL ($\chi^2$ 20$\to$5.8)\\
B1 & ladder at $10^{9}$--$10^{16}$ GeV & stable\\
B2 & up rung with $\kappa$ spread & broad (5--14)\\
\end{tabular}\end{ruledtabular}
\end{table}

\section{Predictions frozen for future data}
\textbf{L1} (equal-coefficient limit, weak): $g_u/g_d=(1+\tfrac1{16})/(m_u+\tfrac1{16})$ with $m_u=8$, i.e.\ $0.1318$; known to be scale-dependent and only meaningful with a model fixing $\kappa$.
\textbf{L2}: no quark family obeys Koide; no lepton family follows the slip scaling.
\textbf{L3}: any new modular fermion family on the $\rho$ ladder has $m_2/m_3\in[\epsilon_m/3,3\epsilon_m]$ for integer $m$; rungs 1 and $\sim\!8$ are occupied.
\textbf{L4}: $\bar\theta=0$ exactly; a $\theta$-induced neutron EDM would falsify the construction.
As a consistency check, LHCb's $B^0\to\mu^+\mu^-$ evidence~\cite{LHCb} agrees with the Standard Model ($R_{\mu\mu}$ $+0.6\sigma$); the Weave adds nothing at this order, so this is consistency, not evidence.

\section{Limits}
The rung formula is computed, but its normalisation and the ``$1/16$'' are the numbers of a specific supergravity class; the coupling constants are free $\mathcal O(1)$ numbers; the up-rung assignment is broad; the quark masses were known before every test except the frozen predictions. Nothing here selects the integers $m$, or explains why the weights are anomaly-free. The result is a map: where quarks can live in the Weave, and how close a zero-extra-parameter version comes.

\textit{Data and code.} \url{https://github.com/dannyeldridge54/AEGIS} (branch \texttt{unify-bg}): \texttt{theory/quarks/}, \texttt{field\_tests/rules/PREREG\_Q*}, \texttt{field\_tests/results/q*}, ladder data sheet \texttt{theory/quarks/LADDER\_DATASHEET.csv}.

\begin{thebibliography}{99}
\bibitem{P1} D.~Eldridge, ``Inflation from the running of the species scale: an $SL(2,\mathbb Z)$ plateau along the string coupling'', Zenodo, doi:10.5281/zenodo.24377930 (2026).
\bibitem{P2} D.~Eldridge, ``Flavour at the end of inflation: Koide's form, trimaximal mixing and a quantised $\bar\theta$'', Zenodo, doi:10.5281/zenodo.24379357 (2026).
\bibitem{FST} F.~Feruglio, A.~Strumia, A.~Titov, JHEP 07 (2023) 027, arXiv:2305.08908; ``Solving the strong CP problem without axions'', arXiv:2406.01689.
\bibitem{NPP} P.~P.~Novichkov, J.~T.~Penedo, S.~T.~Petcov, JHEP 03 (2022) 149, arXiv:2201.02020.
\bibitem{PT} S.~T.~Petcov, M.~Tanimoto, arXiv:2212.13336.
\bibitem{MZ} Quark masses at $M_Z$ (PDG 2024 run), as tabulated in arXiv:2512.13630.
\bibitem{RUN} Running quark and lepton parameters, arXiv:2510.01312.
\bibitem{LHCb} LHCb Collaboration, arXiv:2611.00230.
\bibitem{F17} F.~Feruglio, arXiv:1706.08749.
\end{thebibliography}
\end{document}
